\section{System Overview}
\label{sec:overview}

\subsection{Paper Boundary}

The system is a drone-centered voice safety layer. It listens to short onboard audio windows, converts them into time-frequency features, predicts one of \emph{emergency}, \emph{movement}, or \emph{unknown}, and lets downstream safety logic decide whether a command is safe to dispatch. The paper should not present this as open-ended conversation or as a completed cross-language system. Cross-language variation is motivation for a robust abstraction, not the main contribution in the current evidence set.

\subsection{Current Evidence Stack}

The current repository supports three distinct evidence layers:
\begin{itemize}
\item \textbf{Main model anchor:} \texttt{w14/preprocess\_ext}, using \texttt{base} teacher and student profiles with the clean teacher reused from the \texttt{w14} baseline path.
\item \textbf{Deployment candidate branch:} \texttt{w17/B\_small\_teacher\_student}, using \texttt{xiao\_bottleneck256\_tflm} for both teacher and student profiles.
\item \textbf{Embedded feasibility diagnostics:} Phase 2 ESP32/TFLM artifacts, which expose TFLite operator counts, tensor sizes, arena pressure, and the grouped-convolution incompatibility that motivates the TFLM-specific profile.
\end{itemize}

These layers should not be merged into a single claim. The \texttt{w14} line supports the primary noisy-set performance story. The \texttt{w17} line supports an ongoing compression and compatibility story. The Phase 2 artifacts support a deployment-risk story.

\subsection{End-to-End Flow}

The intended flow is:
\[
\text{Audio window} \rightarrow \text{log-mel frontend} \rightarrow \text{intent model} \rightarrow \text{safety decision} \rightarrow \text{drone command path}.
\]
The current completed benchmark evidence stops at intent prediction. The embedded artifacts show how the selected model would be exported and constrained for ESP32/TFLM, but they do not yet close the full onboard loop.

\subsection{What Is Stable and What Is Ongoing}

The stable anchor is \texttt{w14/preprocess\_ext}: its noisy classification report gives $0.88$ accuracy, $0.79$ emergency recall, and $0.87$ emergency F1 over support $9984$. The closest current TFLM-oriented run, \texttt{w17/B\_small\_teacher\_student}, reports $0.87$ accuracy, $0.79$ emergency recall, and $0.86$ emergency F1 over the same support. That is close enough to motivate continued deployment work, but not enough to replace the anchor.

The ongoing branch is \texttt{xiao\_bottleneck256\_tflm}. The handoff records a target Branchformer entry of $(32,256)$ and $712067$ parameters, with local export sanity checks showing \texttt{CONV\_2D=6} and \texttt{DEPTHWISE\_CONV\_2D=1}. The correct paper wording is therefore ``deployment candidate under active convergence,'' not ``deployed system.''
